\section{Pure Reasoning}

课程里讨论的下一个问题是：能不能把 reasoning 从 knowledge 里尽量分离出来？

\subsection{ARC-AGI}

ARC-AGI 试图测试更接近抽象推理的能力。它的吸引力在于：题目不太依赖背诵事实，而更依赖模式发现、概括和规则迁移。

\begin{itemize}
    \item ARC-AGI-1 已经很难。
    \item ARC-AGI-2 更难，说明``纯推理''并不是一个容易孤立出来的维度。
\end{itemize}

\subsection{MATH 和 AIME}

MATH 数据集包含约 5000 道竞赛级数学题，从 1 级到 5 级难度。AIME（American Invitational Mathematics Examination）则是美国数学竞赛的高水平题目，常被用来测试推理模型的上限。

评估数学题的关键技术细节：
\begin{itemize}
    \item \textbf{答案归一化}：数学题的正确答案可以有多种等价形式（如 $\frac{1}{2}$、$0.5$、$50\%$），评估时需要做符号级别的等价判断，而非简单的字符串匹配。
    \item \textbf{pass@k vs majority vote}：对于同一道题运行 $k$ 次，pass@k 检查是否至少有一次答对，而 majority vote 取最多出现的答案。两种方式给出的``能力''含义不同——前者是``能不能做到''，后者是``稳不稳定''。
    \item \textbf{推理链的正确性}：模型可能给出正确答案但推理过程是错的（碰巧猜对），也可能推理过程合理但在最后一步算错。仅看最终答案的 metric 无法区分这两种情况。
\end{itemize}

\begin{knowledgebox}{推理和知识并不完全可分}
现实模型的表现往往同时依赖知识、语言理解、注意力控制和搜索能力。所谓”纯推理”更多是一种评估理想，而不是一个完全纯净的实验条件。
\end{knowledgebox}

\subsection{本章小结}

纯推理评估面临的核心困境是：推理能力很难完全与知识分离。ARC-AGI 类任务试图最小化知识依赖，但模型在这类任务上的进步往往依赖大量计算（搜索）和对模式的隐式``记忆''。MATH 类任务则需要数学知识和推理能力的结合，两者几乎不可能完全分开。

